%%%%%%%%%%%%%%%%%%%%%%%%%%%%%%%%%%%%%%%%%%%%%%%%%
\subsection{Evaluating Synthetic Image Quality}\label{sec:lit_synthetic_evaluation}

Assessing the quality of generated imagery numerically has historically relied on statistics computed from a pretrained image classifier's feature space rather than on pixel-level comparison to a reference. The \textit{Inception Score} (IS) \cite{salimansImprovedTechniquesTraining2016} exemplifies this idea in its earliest form: propagating generated images through an Inception-v3 model pretrained on ImageNet, it rewards images whose per-image predicted label distribution $p(y|x)$ is low-entropy (confidently assigned to one class) while requiring the marginal distribution over labels $p(y)$ across the whole sample to be high-entropy (diverse across classes), combined as $\mathrm{IS} = \exp(\mathbb{E}_x[\mathrm{KL}(p(y|x) \| p(y))])$ and averaged over at least $50{,}000$ samples. IS was introduced specifically to avoid the cost and inconsistency of paid human annotators and was reported to correlate with human judgment on CIFAR-10 (Mechanical Turk discrimination accuracy dropped from $78.7\%$ to $71.4\%$ once the sample was filtered to the top $1\%$ by IS), but the same paper's own ablations already flag a specific weakness: directly optimizing IS produces adversarial-example-like artifacts rather than genuine quality gains, and by construction IS never compares a generated sample against any statistic of the real data distribution at all.

\textit{Fréchet Inception Distance} (FID) \cite{heuselGANsTrainedTwo2018} closes exactly this gap: real and generated images are both propagated through the same Inception-v3 pool3 activations, each set of activations is modeled as a multivariate Gaussian, and FID is the Fréchet (Wasserstein-2) distance between the two Gaussians, $\mathrm{FID} = \|\mu_r - \mu_g\|^2 + \mathrm{Tr}(C_r + C_g - 2(C_r C_g)^{1/2})$. Introduced alongside a two-time-scale update rule for adversarial training, FID was shown to increase monotonically under six independently controlled disturbance types (Gaussian noise, Gaussian blur, implanted occlusion rectangles, image swirling, salt-and-pepper noise, and cross-dataset contamination), a consistency property IS does not reliably exhibit. \textit{Kernel Inception Distance} (KID) \cite{binkowskiDemystifyingMMDGANs2021} was proposed as a direct fix for FID's principal statistical weakness: FID's Gaussian-moment estimator is provably biased and requires a large sample count before it converges to a stable value (bias was still shrinking at $n=10{,}000$ CIFAR-10 test images in the KID paper's own replication), and the bias is not guaranteed to be consistent across different pairs of distributions, meaning two FID scores computed at the same $n$ can rank two models in the wrong order. KID instead computes the (unbiased, no-Gaussian-assumption) squared Maximum Mean Discrepancy between Inception representations under a cubic polynomial kernel, converges to its asymptotic value even at small $n$, and shares FID's monotonic sensitivity to the same synthetic disturbances.

Automatic feature-space metrics, however, do not always track human perceptual judgment, and human evaluation protocols for generative models were, prior to a deliberate standardization effort, largely ad hoc: informal crowdsourced real/fake discrimination tasks run without a validated design or any check on whether independent raters even agree. \textit{HYPE} \cite{zhouHYPEBenchmarkHuman2019} was introduced as a structured, blind, multi-rater alternative grounded in psychophysics: raters see a balanced, randomized, model-identity-hidden stream of real and generated images -- an explicitly blind design -- and either (HYPE$_{\mathrm{time}}$) have their exposure time adaptively staircased toward an individual perceptual threshold, or (HYPE$_\infty$, the cheaper variant used for most reported comparisons) classify $100$ images ($50$ real, $50$ fake) under no time limit, with the resulting error rate reported alongside bootstrapped $95\%$ confidence intervals ($30$ evaluators, $10{,}000$ resampling iterations) and ANOVA/Tukey tests confirming that models are statistically separable. Applied to \textit{StyleGAN}, \textit{ProGAN}, \textit{BEGAN}, and \textit{WGAN-GP} on CelebA-$64$, even the best model (\textit{StyleGAN} with the truncation trick) deceived only $50.7\%$ of raters -- barely above the $50\%$ chance baseline that defines indistinguishability -- illustrating how far a strong contemporary generator remained from photorealism once measured rigorously rather than by visual inspection. Critically, HYPE's own results provide direct evidence that automatic and human evaluation axes can diverge and must therefore be triangulated rather than substituted for one another: HYPE scores showed no significant correlation with FID on face generation (Spearman $\rho=-0.029$, $p=0.96$) and only weak, model-dependent correlation with FID and KID on ImageNet-5 conditional generation.

Beyond realism, an image can be photorealistic yet fail to depict what its generating prompt asked for -- prompt adherence is a distinct failure mode that IS, FID, and KID cannot detect at all, since none of them examine the conditioning text. \textit{CLIPScore} \cite{hesselCLIPScoreReferencefreeEvaluation2021} closes this gap with a reference-free measure of image-text compatibility: passing an image and a candidate caption or prompt through CLIP's respective encoders and rescaling their cosine similarity, $\mathrm{CLIPScore}(I,C) = 2.5 \cdot \max(\cos(E_I, E_C), 0)$. Despite requiring no ground-truth reference captions at all -- unlike the $n$-gram-overlap metrics (BLEU, CIDEr, SPICE) it was benchmarked against -- CLIPScore achieved the highest correlation with human caption-quality judgments on Flickr8K-Expert ($\tau_c=51.2$ versus $44.9$ for SPICE, the strongest reference-based competitor), demonstrating that a general-purpose vision-language embedding trained on nothing more than web image-caption pairs transfers directly to quality assessment without task-specific supervision.

None of the preceding axes, however, measures what ultimately matters for a thesis whose object is a downstream detector: whether a model trained on the generated imagery performs well on real data. \textit{Classification Accuracy Score} (CAS) \cite{ravuriClassificationAccuracyScore2019} operationalizes this train-on-synthetic, test-on-real idea directly: a classifier is trained exclusively on class-conditional samples drawn from the generative model and then evaluated on the real held-out test set, with real-test top-1/top-5 accuracy serving as the score. Applied to \textit{BigGAN-deep} -- a GAN whose IS and FID scores at the time approached those of real ImageNet data -- CAS revealed a $27.9$-point top-1 and $41.6$-point top-5 accuracy gap relative to training on real data, a degradation invisible to IS and FID alone; conversely, likelihood-based models (\textit{VQ-VAE-2}, Hierarchical Autoregressive Models) with markedly worse FID and IS scores than \textit{BigGAN-deep} substantially outperformed it on CAS. Under a controlled truncation-trick sweep the paper reports essentially no useful correlation between CAS and either IS ($\rho=-0.86$, the wrong sign) or FID ($\rho=0.16$), and CAS additionally exposed specific classes (balloon, paddlewheel, pencil sharpener, spatula) on which \textit{BigGAN-deep}'s downstream classifier scored $0\%$ on the real validation set -- failures a single global FID number cannot localize.

Taken together, these three axes -- structured blind human rating, automatic distributional and embedding-based proxies, and downstream task utility -- are exactly the three-way triangulation this thesis adopts when evaluating its own candidate synthetic-image generators (\Cref{sec:results_synthetic_data_practicability}), following the evaluation methodology in \texttt{docs/synthetic-model-comparison/06\_evaluation-methodology.md}: a structured, blind, multi-rater qualitative rubric in the spirit of HYPE's design; automatic proxies drawing on FID and KID for distributional realism, CLIPScore for prompt adherence, and a teacher-classifier-confidence-on-synthetic score that isolates exactly the per-image confidence half of the Inception Score's original formulation \cite{salimansImprovedTechniquesTraining2016} -- without its population-diversity term, since the target here is recognizability by a fixed downstream teacher rather than raw sample diversity -- as a cheap, high-sample-count recognizability proxy; and a CAS-style downstream axis, reported as real-test detector mAP after training on a given generator's synthetic images, following the general train-on-synthetic-test-on-real pattern \cite{ravuriClassificationAccuracyScore2019}. No single axis is treated as decisive precisely because HYPE and CAS both demonstrate empirically that these axes can disagree: a generator strong on one need not be strong on another, and only agreement across all three constitutes a robust claim about a generator's fitness for this thesis's actual downstream use.

%%%%%%%%%%%%%%%%%%%%%%%%%%%%%%%%%%%%%%%%%%%%%%%%%
\subsection{Synthetic Data for Long-Tail and Rare-Class Problems}\label{sec:lit_synthetic_long_tail}

Real-world visual recognition datasets rarely present a uniform class distribution: a small number of classes dominate the sample count while most classes are supported by comparatively few examples, a regime the long-tailed learning literature groups into two broad families of remedy, \textit{class re-balancing} (re-sampling and cost-sensitive re-weighting) and \textit{information augmentation} (transfer learning and data augmentation) \cite{zhangDeepLongTailedLearning2023}. Both families predate the current generation of generative image models and were originally developed to make better use of a class's \textit{existing} samples rather than to manufacture new ones.

\textit{SMOTE} \cite{chawlaSMOTESyntheticMinority2002} is the foundational instance of the re-sampling family and a useful contrast case for this thesis's approach: rather than duplicating minority-class examples verbatim (which was already known to cause overfitting to a narrower, more specific decision region), it synthesizes new minority-class points along the line segment joining a minority sample to one of its $k=5$ nearest minority-class neighbors \textit{in feature space}, evaluated via ROC/AUC gains on tabular and pixel-feature classifiers (C4.5, Ripper, Naive Bayes) rather than on raw images. This is a categorical limitation for the present use case: interpolating between existing feature vectors can only recombine visual information the dataset already contains (the same pose, background, and lighting variation already present in the sampled neighbors); it cannot supply a camera angle, habitat, or lighting condition that no real photograph of the class exhibits, which is precisely the kind of novel appearance content raw generative image synthesis is used to add in \Cref{sec:synthetic_generation_pipeline}'s six axes of variation.

Cost-sensitive re-weighting takes a different, loss-side approach: rather than changing what data the model sees, it changes how much each sample's error counts. \textit{Class-Balanced Loss} \cite{cuiClassBalancedLossBased2019} reweights any per-sample loss (softmax cross-entropy, sigmoid cross-entropy, or focal loss) by the inverse of a class's \textit{effective number of samples}, $E_n = (1-\beta^n)/(1-\beta)$, a quantity derived from a random-covering argument that captures the diminishing marginal information contributed by additional, likely-overlapping samples of an already well-represented class; the hyperparameter $\beta \in [0,1)$ smoothly interpolates between no re-weighting ($\beta=0$) and classical inverse-class-frequency re-weighting ($\beta \to 1$), and the method was validated on long-tailed CIFAR, ImageNet, and \textit{iNaturalist}. \textit{LDAM} \cite{caoLearningImbalancedDatasets2019} instead re-margins rather than re-weights: starting from a per-class generalization-error bound, it derives an optimal per-class margin trade-off ($\gamma_j \propto n_j^{-1/4}$ for class $j$ with $n_j$ training samples) and encourages minority classes to maintain a proportionally larger decision-boundary margin, combined with a \textit{deferred} re-weighting schedule (training first with the unweighted loss, re-weighting only in a later stage) to sidestep the optimization difficulties that re-weighting from the start of training is known to cause. Both methods are explicit about operating purely on the training objective, applied on top of whatever real images a class already has -- they change how existing samples are counted, not whether enough of them exist to represent the class's visual diversity in the first place.

A third line of work questions whether re-balancing the classifier is even the right target at all. \textit{Decoupling} \cite{kangDecouplingRepresentationClassifier2020} separates representation learning from classifier learning and finds, across ImageNet-LT, Places-LT, and \textit{iNaturalist}, that the best feature representations are learned under plain instance-balanced sampling (i.e., no re-balancing at all during representation learning), and that long-tailed recognition performance is instead recovered by re-adjusting only the classifier afterward -- retraining it with class-balanced sampling, replacing it with a nearest-class-mean classifier, or simply renormalizing its per-class weight norms ($\tau$-normalization) -- consistently outperforming methods that re-balance sampling or the loss throughout training. This result reframes the practical failure mode of long-tailed recognition as primarily a biased \textit{decision boundary} rather than a poorly learned \textit{representation}, but it presupposes that a class's representation was learnable from its available images to begin with; it offers no remedy for a class whose real-image pool is so small that no representation, however sampled, can be learned from it.

This is exactly the gap that \citeauthor{zhangDeepLongTailedLearning2023}'s taxonomy of information-augmentation methods \cite{zhangDeepLongTailedLearning2023} leaves open. Within its data-augmentation sub-category, the survey separates \textit{head-to-tail transfer augmentation} -- translating or perturbing head-class samples toward a tail class (M2m) or displacing tail-class feature vectors along directions estimated from head-class feature statistics (RSG, FTL, LEAP) -- from \textit{non-transfer augmentation}, which includes deep-learning adaptations of SMOTE-style feature interpolation, data mixup (MiSLAS, Remix), and Gaussian-prior feature synthesis (FASA, MetaSAug). Every one of these methods, transfer-based or not, operates in \textit{feature space} on top of samples the dataset already contains; none synthesizes new raw images from an external generative prior. The survey's own discussion is explicit that this remains unresolved: \enquote{how to better conduct data augmentation for long-tailed learning is still an open question,} and separately warns that naively applying class-agnostic augmentation to a long-tailed dataset augments head classes as much as tail classes, providing no relative benefit to the classes that need it most.

Taken together, this literature motivates rather than substitutes for the band structure adopted in this thesis (\Cref{sec:band_structure}). Re-sampling, cost-sensitive re-weighting, and feature-space augmentation all share a structural limitation for the most data-poor classes: they rearrange, reweight, or interpolate information that a class's existing image pool already contains, so their ceiling is set by how much visual diversity that pool holds in the first place -- SMOTE-style interpolation cannot exceed the feature-space span of its neighbors, class-balanced loss and LDAM only change how existing images are counted during training, and decoupling's classifier-only fix presupposes a representation was learnable from the class's images to begin with. Image-level generative synthesis is a categorically different intervention: rather than adjusting how existing samples are counted, weighted, or interpolated, it manufactures new appearance content -- camera angle, habitat, lighting, pose -- that a data-poor class's own real-image pool cannot supply. This is precisely why synthetic supplementation in this thesis is reserved for Band A and Band B, the two bands whose post-curation real-image pool ($<150$ and $150$--$249$ images respectively) is judged too small to support adequate visual diversity on its own, while Band C reaches the same $200$-image training budget entirely from real imagery, deliberately withholding synthetic supplementation so it can serve as a real-data control against the synthetic-heavy bands at matched training-set size.
